\begin{abstract}
Agentic workflows allow complex tasks to be carried from a natural language prompt. This has
broad applicability to allow non-expert users to perform difficult work. To facilitate the
agents' ability to do things beyond text output, tools need to be provided in the form of an
API. We examine two ways to expose this interface: A model context protocol (MCP) server which
shows all available processes along with metadata, and a vector database populated with the
available processes and accessed using retrieval augmented generation (RAG). These two
approaches are compared within an agentic system to create a workflow for a sample geospatial
analysis task, and the input prompt is paraphrased in various ways to improve the robustness of
the trial. Paraphrases are systematically created using an LLM-generated and human-reviewed
strategy. The resulting workflows are evaluated manually and processed with the actual tools to
ensure correctness. Our findings indicate that MCP provides higher accuracy and greater
reliability when compared to RAG. Additionally, paraphrasing did not cause a substantial
difference in performance, showing that the agentic system can effectively process a range of
distinct inputs.
\end{abstract}
